\subsection{A concrete positive-trace budget from the original period kernels}
\label{sec:retained-period-window-budget}

This statement isolates the remaining quantitative inputs on the
actual period family.  The normalizers are the original seed
normalizers, and the total calibration is prescribed independently
of the observed trace.  The hypotheses below are estimates for
the literal period Gram blocks, not for substitute column maps.

Let \(A_u=c_{0X,u}^{-1/2}P_u:C_u\to H\) be the original
finite family, with \(d_u=\dim C_u\), \(D=\dim H>0\),
\(R=\sum_u d_u>0\), and \(D_*>0\) the original proposed
dimension main term.  Put
\[
 K_{uu'}=A_u^*A_{u'},\qquad
 S=\bigoplus_uK_{uu},\qquad Q=\sum_uA_uA_u^*.
                                                        \tag{PW1}
\]
Assume the actual kernel estimates
\[
 \|K_{uu}-I\|_{\rm HS}^2\le q^2d_u,
 \qquad
 \|K_{uu'}\|_{\rm HS}^2
       \le(1+\tau)\frac{d_ud_{u'}}{D_*}\quad(u\ne u'),
 \qquad q,\tau\ge0 .                                  \tag{PW2}
\]
These are the diagonal defect and the distinct-word estimate.
In the original source route one may take
\[
 q=\sqrt{x_K}+\sqrt{C_0/\Gamma},\qquad
 \tau=b^{-6},\qquad x_K=\exp(-kf\log k),                 \tag{PW3}
\]
once the nonzero-Gram diagonal remainder and the distinct-word
scalar estimate have been established.  Formula (PW3) is a
consequence of those estimates, not a replacement for them.

\begin{theorem}\label{thm:retained-period-window-budget}
Write
\[
 \sigma=D_*/R,\qquad \xi=|D/D_*-1|,\qquad
 \mathcal E=\frac12\left(q+
  \sqrt{(1+\xi)\bigl(1+\tau+(1+q)^2\sigma\bigr)-1+2q}
                         \right).                      \tag{PW4}
\]
At the prescribed calibration \(c=R\), the original source
and target positive traces satisfy \(g_R+b_R\le\mathcal E\).
For actual source isometries \(j_u:C_u^0\to C_u\) and an
actual target projection \(\Pi\), write
\[
 R_0=\sum_u\dim C_u^0>0,\qquad d=\operatorname{rank}\Pi>0,
 \qquad r_0=R_0/R,\qquad v_0=d/D.
\]
If the prescribed retained calibration is \(c_0=R_0\), then
the same retained maps \(\Pi A_uj_u\) satisfy
\[
 g_0+b_0\le
       \frac{\mathcal E+\max\{1-r_0,1-v_0\}}{r_0}.
                                                        \tag{PW5}
\]
For an arbitrary prescribed \(c_0=\theta R>0\), the conclusion
instead retains the exact calibration terms:
\[
 g_0+b_0\le
 \frac{\mathcal E+1-v_0+(\theta-r_0)_++(v_0-\theta)_+}
      {\theta}.                                       \tag{PW6}
\]
No equality between \(c_0\) and an observed trace is used.
\end{theorem}

\begin{proof}
Let \(M=\operatorname{tr}S=\operatorname{tr}Q\).  Summing
the diagonal estimate in (PW2) gives
\[
 \|S-I\|_{\rm HS}^2\le q^2R,
       \qquad |M-R|\le qR.                             \tag{PW7}
\]
The second assertion is Cauchy--Schwarz on the whole direct
sum.  Also
\(\|K_{uu}\|_{\rm HS}\le(1+q)\sqrt{d_u}\).
Using the same common target for every column gives the exact
identity
\[
 \operatorname{tr}(Q^2)=\sum_{u,u'}\|K_{uu'}\|_{\rm HS}^2.
\]
Consequently
\[
 \operatorname{tr}(Q^2)
       \le(1+\tau)R^2/D_*+(1+q)^2R.                    \tag{PW8}
\]
The discarded diagonal subtraction in the off-diagonal sum
is nonnegative.  Equations (PW7)--(PW8) imply
\[
 \frac{D\operatorname{tr}(Q^2)-2RM+R^2}{R^2}
 \le(1+\xi)\bigl(1+\tau+(1+q)^2\sigma\bigr)-1+2q.
                                                        \tag{PW9}
\]
The right-hand side is nonnegative because all four parameters
\(\xi,\tau,\sigma,q\) are nonnegative.  The exact prescribed
trace identity and the Hilbert--Schmidt bound in
Theorem~\ref{thm:retained-trace-fixed-moment}, applied with
\(c=R\), now give (PW4).  Apply the joint source/target
restriction inequality with this same prescribed calibration
to obtain (PW6).  For \(\theta=r_0\), its extra numerator is
\[
       1-v_0+(v_0-r_0)_+=\max\{1-v_0,1-r_0\},
\]
which gives (PW5).
\end{proof}

For example, when \(q,\xi\le1\),
\[
 \mathcal E\le\tfrac12q+
                  \tfrac12\sqrt{\xi+2\tau+8\sigma+2q}.
                                                        \tag{PW10}
\]
Thus quantitative bounds for the four original numbers
\(q,\xi,\tau,\sigma\), together with the actual retained
rank fractions and calibration, give a numerical bound at
every requested window.  At each arithmetic parameter this
may be substituted into the established norm penalty as
\[
 P_{\rm norm}(x)\le\kappa(x)^{-1}
       \operatorname{avg}_u
           \min\{H_u,\mathcal E_{0,u}(U_u+H_u)\},        \tag{PW11}
\]
where \(\mathcal E_{0,u}\) is (PW5) or (PW6) for the actual
maps and scalars.  To prove the moving arithmetic budget one
must establish those bounds uniformly at the required samples
and identify the retained theta columns with these normalized
period columns, including every scalar modulus.  The newly
proved pointwise estimate and weighted point sum alone do not
supply the distinct-word leading constant \(1+\tau\) in (PW2),
the dimension comparison \(\xi\), or that common-row identity.
